% Chapter 01 — Introduction: Why eBPF?
% Yêu cầu từ main: \ChapterDivider{...}{...}{...}; không cần khai báo lại macro.
\section{Chương 1: Giới thiệu eBPF}
\ChapterDivider{Chương 1}{Giới thiệu eBPF}{Từ nhu cầu quan sát Linux đến công nghệ eBPF}

\begin{frame}{Vì sao cần quan sát hoạt động bên trong Linux?}
  \begin{block}{Một tình huống quen thuộc}
    Một máy chủ Linux chạy nhiều dịch vụ cùng lúc. Khi có sự cố, log ứng dụng
    không phải lúc nào cũng cho biết điều gì diễn ra trong kernel.
  \end{block}
  \vspace{0.13cm}
  \begin{columns}[T,onlytextwidth]
    \begin{column}{0.49\textwidth}
      \begin{block}{Ta muốn biết điều gì?}
        \begin{itemize}
          \item Tiến trình nào đang mở tệp hoặc giao tiếp mạng?
          \item Những system call nào được thực hiện?
          \item CPU và I/O đang được sử dụng như thế nào?
        \end{itemize}
      \end{block}
    \end{column}
    \begin{column}{0.49\textwidth}
      \begin{block}{Vấn đề thực tế}
        \centering
        \textbf{Nhiều tiến trình}\\[0.14cm]
        $\downarrow$\\[0.14cm]
        \textbf{Linux kernel}\\[0.14cm]
        \footnotesize{File system \quad Network \quad Scheduler}
      \end{block}
    \end{column}
  \end{columns}
  \vspace{0.10cm}
  \centering\textbf{Câu hỏi:} Làm sao theo dõi các sự kiện này\ mà không cần sửa mã nguồn ứng dụng?
\end{frame}

\begin{frame}{Linux đã có những cách quan sát nào?}
  \begin{columns}[T,onlytextwidth]
    \begin{column}{0.49\textwidth}
      \begin{block}{Application logging}
        Ghi những thông tin mà ứng dụng chủ động cung cấp.
        Muốn thêm một sự kiện vào log có thể phải sửa ứng dụng.
      \end{block}
      \begin{block}{\texttt{strace}}
        Theo dõi các system call của tiến trình; hữu ích khi cần xem
        một chương trình tương tác với kernel ra sao.
      \end{block}
    \end{column}
    \begin{column}{0.49\textwidth}
      \begin{block}{\texttt{perf}}
        Thu thập và phân tích dữ liệu hoạt động, hiệu năng; cách lấy mẫu
        và chi phí quan sát phụ thuộc phép đo.
      \end{block}
      \begin{block}{Kernel module}
        Có thể bổ sung chức năng cho kernel, nhưng mã chạy với đặc quyền cao;
        sai sót có thể ảnh hưởng toàn hệ thống.
      \end{block}
    \end{column}
  \end{columns}
  \vspace{0.02cm}
  \begin{alertblock}{eBPF nằm ở đâu trong bức tranh này?}
    eBPF là \textbf{một lựa chọn bổ sung} để xử lý sự kiện tại các hook được kernel hỗ trợ.
  \end{alertblock}
\end{frame}

\begin{frame}{eBPF là gì?}
  \begin{block}{Khái niệm}
    \textbf{eBPF} cho phép nạp và chạy những chương trình đã được kernel kiểm tra,
    nhằm theo dõi hoặc xử lý sự kiện tại các \textit{hook} tương thích.
  \end{block}
  \vspace{0.16cm}
  \begin{columns}[T,onlytextwidth]
    \begin{column}{0.32\textwidth}
      \begin{block}{Theo sự kiện}
        Một sự kiện phù hợp kích hoạt chương trình eBPF được gắn vào hook.
      \end{block}
    \end{column}
    \begin{column}{0.32\textwidth}
      \begin{block}{Xử lý tại kernel}
        Có thể lọc hoặc tổng hợp dữ liệu trước khi gửi ra User Space.
      \end{block}
    \end{column}
    \begin{column}{0.32\textwidth}
      \begin{block}{Chia sẻ kết quả}
        Ứng dụng ở User Space nhận sự kiện hoặc đọc dữ liệu từ map.
      \end{block}
    \end{column}
  \end{columns}
  \vspace{0.16cm}
  \begin{alertblock}{Điều cần nhớ}
    eBPF không chỉ dùng để lọc gói tin: nó còn phục vụ quan sát hệ thống,
    xử lý mạng và một số cơ chế bảo mật.
  \end{alertblock}
\end{frame}
